\section{Large Language Models}
\label{chap:background:sec:llms}

Large Language Models (LLMs) are generative transformer models trained to perform natural language processing tasks.
State-of-the-art LLMs typically employ a decoder-only transformer architecture \cite{radford2018improving}, consisting of a stack of layers, each containing an attention block \cite{attention_is_all_you_need} and a feed-forward network, both of which are surrounded by normalization layers.
Although this structure is broadly shared, different LLMs have substantial differences in their architecture.
The first key distinctive aspect is the attention mechanism. The original \textit{Attention is all you need} paper \cite{attention_is_all_you_need} proposed Multi-Head Attention (MHA), which comprises multiple attention heads, each performing a three-fold projection of the input into the query $Q$, key $K$, and value $V$ latent spaces. 
This mechanism can be formalized as: 

\[
\text{MultiHead}(Q, K, V) = \text{Concat}\big(\text{head}_1, \ldots, \text{head}_h\big)
\]
\[
\text{where} \quad \text{head}_i = \text{Attention}\big(Q W_i^Q,\, K W_i^K,\, V W_i^V\big)
\]
\[
\text{and} \quad \text{Attention}(Q, K, V) = \text{softmax}\!\left(\frac{Q K^\top}{\sqrt{d_k}}\right) V.
\]
More recent models employ Grouped Query Attention (GQA) \cite{ainslie2023gqa}, which divides the attention heads into groups that share $K$ and $V$ projections, while performing a distinct query $Q$ projection for each group. This allows for reducing the amount of computation and the memory footprint of inference, as $K$ and $V$ projections of previous tokens are typically cached, a technique known as KV caching. Notably, DeepSeek V3 and R1 \cite{liu2024deepseek, guo2025deepseek} use a different mechanism called Multi-head Latent Attention (MLA), which compresses the KV cache for further memory reduction.

The second aspect in which large language models differ is the sparsity of the feed-forward block. Dense models utilize monolithic feed-forward networks, whereas Mixture-of-Experts (MoE) \cite{artetxe2021efficient} models comprise a set of smaller blocks, known as experts, that are selectively activated by a router network. MoE models achieve competitive accuracy compared to similarly large dense models, while necessitating a fraction of the compute.

Finally, models vary widely in terms of size, ranging from hundreds of millions to trillions of parameters. The factors that contribute to the number of parameters $N$ are the number of layers $n_\text{layer}$, attention heads $n_\text{heads}$, and their output size $d_{\text{attn}}$, and the dimension of the feed-forward block $d_{\text{ff}}$.
